% The byte path costs no processor cycles. Three events call one function.
\begin{tikzpicture}[font=\sffamily\small]
\node[hw, text width=24mm] (usart) at (0,1.9) {USART3\\{\scriptsize receive request}};
\node[hw, text width=24mm] (mux) at (3.1,1.9) {request multiplexer\\{\scriptsize number from RM0455}};
\node[master, text width=24mm] (dma) at (6.3,1.9) {DMA1 stream 0\\{\scriptsize circular, NDTR}};
\node[memblk, text width=26mm] (buf) at (9.5,1.9) {rx\_buf[512]\\{\scriptsize main memory, 32 B aligned}};
\begin{scope}[on background layer]
\node[layer, fit=(usart)(mux)(dma)(buf), inner ysep=7pt,
      label={[layerlabel, anchor=south west]north west:the byte path: no processor cycles at all}] (bp) {};
\end{scope}

\node[fw, text width=32mm] (chk) at (3.8,-0.9) {rx\_check()\\{\scriptsize pos = 512 minus NDTR}};
\node[fw, text width=28mm] (inv) at (9.5,-0.9) {invalidate by address\\{\scriptsize argued in chapter 19}};
\node[user, text width=28mm] (cons) at (9.5,-3.2) {consumer\\{\scriptsize reads the span in place}};

\draw[flow] (usart) -- (mux);
\draw[flow] (mux) -- (dma);
\draw[flow] (dma) -- (buf);
\draw[irq] (usart.south) |- node[lbl, above, pos=0.72]{IDLE} (chk.west);
\draw[irq] (dma.south) -- node[lbl, right, xshift=1mm]{HT, TC} (chk.north east);
\draw[flow] (chk) -- node[lbl, above]{pointer, length} (inv);
\draw[bus] (buf.south) -- (inv.north);
\draw[flow] (inv) -- node[lbl, right, xshift=1mm]{clean bytes} (cons);

\node[umlnote, text width=44mm, anchor=north west] at (-2.0,-2.3)
  {Half transfer, transfer complete and line idle are three triggers for one
   function. Three separate code paths would become three copies of the same
   arithmetic, and the least used copy is the one that is wrong.};
\end{tikzpicture}
